%
% The Pierre Auger Observatory (chapter 3)
%
\chapter{The Pierre Auger Observatory}\label{chap:auger}

The Pierre Auger Observatory is located near the town of Malarg\"ue in the Mendoza province of Argentina between the Andes mountain range and la Pampa Amarilla.
Auger, as the observatory is commonly referred to, is the world's largest cosmic-ray observatory, covering a total area of \qty{3000}{\kilo\meter^{2}}.
In context, this is about the size of the state of Rhode Island, or for those with a more worldly perspective, slightly larger than the size of Luxembourg.
The observatory's size allows the detection of ultra-high-energy cosmic rays, starting at energies of $\approx$\qty[parse-numbers=false]{10^{17}}{\electronvolt}, with high statistics.
However, the observatory is ground based and therefore indirectly detects the highest energy particles in the known Universe through the air showers they initiate after colliding with atmospheric nuclei.

Not only is Auger the world's largest cosmic-ray observatory, but adding to the novelty is the hybrid detection mode employed by the observatory.
The observatory consists of two primary detectors, the surface detector (SD), an array of over 1600 water-Cherenkov detectors (WCDs) spread across the full \qty{3000}{\kilo\meter^{2}} footprint of the observatory, and a fluorescence detector (FD), 24 fluorescence telescopes overlooking the SD array at four locations around the edges of the array.
Auger has been recording data in hybrid mode for over 20 years, officially starting in 2004 and with full completion of the SD in 2008.
In March 2026, Auger finished deploying detectors for the next-generation upgrade of the observatory, AugerPrime.
An international funding agreement will keep the observatory in operation until at least 2035, allowing for high-quality air-shower measurements with AugerPrime, and the potential for prototyping future detectors in the Argentinian Pampa.
This chapter will describe the Pierre Auger Observatory and its detectors systems in detail.


\section{Surface Detectors}\label{sec:sd}

\section{Fluorescence Detectors}\label{sec:fd}

\section{Event Reconstruction}\label{sec:reconstruction}

\section{Auger Engineering Radio Array}\label{sec:aera}



\section{AugerPrime: The Upgrade of the Pierre Auger Observatory}\label{sec:augerprime}

\subsection{Surface Scintillator Detector}\label{sec:ssd}


\subsection{Underground Muon Detectors}\label{sec:umd}
